\documentclass[11pt]{article}
\usepackage[utf8]{inputenc}
\usepackage[T1]{fontenc}

\usepackage{amsmath,amssymb,amsthm}
\newtheorem{theorem}{Théorème}
\newcommand{\R}{\mathbb{R}}
\newcommand{\N}{\mathbb{N}}
\DeclareMathOperator{\Ker}{Ker}
\title{Algèbre linéaire}
\author{Poly de cours}
\begin{document}
\maketitle
\section{Applications linéaires}
Soit $f : \R^n \to \R^p$ une application linéaire. Son noyau est $\Ker f = \{x \in \R^n \mid f(x) = 0\}$.
\begin{theorem}[Rang]
Pour tout $f \in \mathcal{L}(\R^n, \R^p)$, on a
\begin{equation}\label{eq:rang}
\dim \Ker f + \operatorname{rg} f = n.
\end{equation}
\end{theorem}
D'après \eqref{eq:rang}, si $n > p$ alors $f$ n'est pas injective.
\begin{align}
\det(AB) &= \det(A)\det(B) \label{eq:det} \\
\det(A^{-1}) &= \frac{1}{\det A} \nonumber
\end{align}
\[
M = \begin{bmatrix} a & b \\ c & d \end{bmatrix}, \qquad
\left| \begin{matrix} a & b \\ c & d \end{matrix} \right| = ad - bc
\]
Pour $n \in \N$, $\displaystyle \sum_{k=0}^{n} \binom{n}{k} = 2^n$.
\end{document}
